% !TeX root = ../main.tex
%%% ---------------------------------------------------------------------------
%%% 附录
%%%
%%% 位置：参考文献**之后**（GB/T 7713）。main.tex 里的 \appendix 已经把
%%% 编号切成「附录 A」，图表公式自动变成 图 A1 / 表 A1 / 式 (A1)。
%%%
%%% 放什么：过长的推导、完整的参数表、补充数据——凡是放正文会打断阅读、
%%% 但审稿人可能需要核对的内容。
%%% ---------------------------------------------------------------------------

\section{方差上界的紧性讨论}
\label{app:tightness}

\cref{thm:variance} 给出的界在 $g \in \setof{0, 1}$ 时取到等号。对
$g \in (0,1)$，由\cref{eq:var-expand} 可以进一步给出

\begin{equation}
  \label{eq:tight}
  \Var\!\left[\featmap_l'\right]
    \ge \frac{\sigma_1^2 \sigma_2^2}{\sigma_1^2 + \sigma_2^2},
\end{equation}

等号在 $g^{*} = \sigma_2^2 / (\sigma_1^2 + \sigma_2^2)$ 处取得。这说明
门控不仅不放大方差，在中间取值处实际上起到了方差收缩的作用。

\section{完整超参数配置}
\label{app:hyperparams}

\begin{table}[htbp]
  \centering
  \caption{完整的训练超参数}
  \label{tab:hyperparams}
  \begin{tabular}{l l}
    \toprule
    {超参数} & {取值} \\
    \midrule
    优化器            & SGD（动量 \num{0.9}） \\
    初始学习率        & \num{0.01} \\
    学习率策略        & 余弦退火，预热 \num{500} 步 \\
    权重衰减          & \num{1e-4} \\
    批大小            & 16（每卡 4） \\
    训练轮数          & \qty{36}{\epoch} \\
    输入尺寸          & $1024 \times 1024$ \\
    数据增强          & 随机水平/垂直翻转、随机旋转 $90^{\circ}$ \\
    门控压缩比 $r$    & 16 \\
    损失权重 $\lambda$ & \num{1.0} \\
    \bottomrule
  \end{tabular}
\end{table}
